\begin{explanationbox}
	\textbf{\textcolor{CExplanation}{一句话直觉}}

	焦点三角形面积只依赖于短半轴$b$和顶角$\theta$，无需计算边长，直接代入半角正切。

	\medskip
	\textbf{\textcolor{CExplanation}{核心拆解}}
	\begin{description}[style=nextline, leftmargin=2.8em, labelsep=0.5em, itemsep=0.45em]
		\item[\textbf{要点一（比例关系）：}] 面积正比于$b^{2}$，比例系数是$\tan\frac{\theta}{2}$。
		\item[\textbf{要点二（半角正切）：}] 使用半角正切可避免因$\theta$锐钝导致的正负讨论。
	\end{description}

	\medskip
	\textbf{\textcolor{CExplanation}{几何本质（定和与余弦）}}

	椭圆上任意一点到两焦点的距离之和为$2a$，通过余弦定理可建立$|PF_{1}||PF_{2}|$与$\theta$的关系，面积由$\frac{1}{2}mn\sin\theta$导出，最终出现半角正切。

	\medskip
	\textbf{\textcolor{CExplanation}{代数意义（三角变换）}}

	公式融合了椭圆定义、余弦定理和半角恒等式，简化了面积计算，体现了$b$的主导作用。

	\medskip
	\textbf{\textcolor{CExplanation}{考点价值（速算应用）}}
	\begin{itemize}[leftmargin=*, itemsep=0.5em, topsep=0.4em]
		\item \textbf{考法一（直接求面积）：} 给出椭圆方程和$\theta$，直接用$S=b^{2}\tan\frac{\theta}{2}$。
		\item \textbf{考法二（逆向求参数）：} 已知面积和椭圆方程，通过$\tan\frac{\theta}{2}=\frac{S}{b^{2}}$求$\theta$或椭圆参数。
	\end{itemize}

	\medskip
	\textbf{\textcolor{CExplanation}{顿悟点}}

	面积公式不含$a$，说明焦点三角形面积与长半轴无关，只由$b$和$\theta$决定，这是椭圆几何性质的直观体现。

	\medskip
	\textbf{\textcolor{CExplanation}{使用场景}}
	\begin{itemize}[leftmargin=*, itemsep=0.5em, topsep=0.4em]
		\item \textbf{场景一（焦点三角形面积）：} 已知椭圆方程和$\angle F_{1}PF_{2}$，快速求面积。
		\item \textbf{场景二（极值问题）：} 结合$\tan\frac{\theta}{2}$的范围可研究面积最值。
	\end{itemize}
\end{explanationbox}
